% ---------------------------------------------------------------------------------------------------------------------
\subsection{Study~II: Multi-Stream Classifier Architecture and Capacity-Matched Controls}
\label{sec:m-study2}

\subsubsection{Rationale}
With the frozen suppressor ensemble $\bar f$ of Study~I, a classifier can receive the radiograph $X$, the soft-tissue image
$S=X-\bar f(X)$ and the bone image $\hat B=\mathrm{clip}_{[0,1]}(X-S)$ on the stored 8-bit scale. Because $(S,\hat B)$ is a deterministic
function of $X$, the inputs carry no new information about a label $Y$:
\begin{equation}
I\big(Y;\,X,S,\hat B\big)=I\big(Y;\,X\big).
\label{eq:info}
\end{equation}
A gain can therefore come only from the prior the suppressor learned from CT, or from the inductive bias of the network that reads the
extra images; a network that reads three images also spends three times the encoder computation. The study separates these sources with a
control of the same capacity, and it analyses the effect per finding group, because suppression could help focal lesions behind ribs while
removing cues for diffuse or pleural findings.

\subsubsection{Input configurations and network}
Five configurations were trained with identical settings (Figs.~\ref{fig:architecture}b and~\ref{fig:cls_design}):
\emph{orig} ($X$), \emph{soft} ($S$), \emph{early} ($X$, $S$ and $\hat B$ as the three input channels of one stream), \emph{hybrid} (three
streams $X$, $S$, $\hat B$) and \emph{ctrl} (the hybrid network with $X$ in all three streams). The soft-tissue image comes from our
suppressor ensemble run at its input size of 512\,px; classifiers see all images at $256\times256$\,px. One DenseNet-121 encoder $E$
\cite{huang2017densenet} with shared weights, pretrained on ImageNet \cite{deng2009imagenet,wightman2019timm}, processes each stream $u_s$.
Each stream is standardised per image and passed through its own learned affine map, because the bone image spans a narrow intensity range;
the $8\times8$ encoder output is projected to $D=256$-dimensional tokens with a stream embedding $e_s$ and a two-dimensional sine position
code $\pi_n$:
\begin{equation}
\begin{gathered}
\tilde u_s=\gamma_s\,\frac{u_s-\mathrm{mean}(u_s)}{\mathrm{sd}(u_s)}+\beta_s,\\
z_{s,n}=W\,E(\tilde u_s)_n+e_s+\pi_n\in\mathbb{R}^{D}.
\end{gathered}
\label{eq:tokens}
\end{equation}
A two-layer decoder with eight heads and one learned query $q_c$ per class attends to the tokens of all streams, as in query-based
multi-label heads \cite{liu2021query2label,ridnik2023mldecoder}, and a linear map per class gives the logit $\ell_c$ from the decoder
output $o_c$. With keys $k_{s,n}$, the head-averaged cross-attention of the last layer defines the routing of class $c$ to stream $s$:
\begin{equation}
\begin{gathered}
a_{c,(s,n)}=\operatorname*{softmax}_{(s,n)}\Big(\frac{q_c^{\top}k_{s,n}}{\sqrt{D}}\Big),\\
\rho_{c,s}=\sum_{n}a_{c,(s,n)},\qquad
\ell_c=w_c^{\top}o_c+b_c.
\end{gathered}
\label{eq:routing}
\end{equation}
Single-stream configurations use the same encoder and decoder over 64 tokens. The shared encoder gives the hybrid the parameters of a
single-stream model plus the decoder, at three times the encoder computation; \emph{ctrl} has exactly the hybrid's parameters and
computation but no suppression.

\begin{figure*}[!t]
\centering
\safefig[0.95\textwidth]{figures/fig_study_design.pdf}
\caption{Classifier study design. (a) The three inputs for an NIH ChestX-ray14 test image: the radiograph $X$, the soft-tissue image $S$
predicted by the suppressor and the bone image $\hat B=X-S$ (box: expert annotation of a focal finding). (b) The five configurations, all
with the same DenseNet-121 encoder, class-query decoder and training settings. In the hybrid and the control one encoder with shared
weights processes each stream and the decoder attends to the tokens of all three; the control receives $X$ three times. Bottom: the
pre-specified contrasts.}
\label{fig:cls_design}
\end{figure*}

\subsubsection{Training}
The loss is the binary cross-entropy over the annotated labels $\mathcal{V}$ of a batch (one-hot targets for the three BDCXR-3257 classes):
\begin{equation}
\mathcal{L}_{\mathrm{cls}}=-\frac{1}{|\mathcal{V}|}\sum_{(i,c)\in\mathcal{V}}\Big[y_{ic}\log\sigma(\ell_{ic})+(1-y_{ic})\log\big(1-\sigma(\ell_{ic})\big)\Big].
\label{eq:bce}
\end{equation}
Training used AdamW \cite{loshchilov2019adamw} (learning rate $3\times10^{-4}$, weight decay $10^{-4}$), batch size 32, 20 epochs with a
cosine schedule after 200 warm-up steps, and mixed precision. Augmentation: random rotation by up to $7^\circ$, scaling by 0.9--1.1,
translation by up to 3\,\% of the image size, and brightness and contrast changes, identical across the channels of an image. The epoch with
the best validation macro-AUROC was kept. Every configuration was trained with seeds 0--4 on NIH ChestX-ray14 and VinBigData and with seeds
0--4 in each of the five BDCXR-3257 folds of Section~\ref{sec:m-bdcxr}, giving 175 runs. A configuration's prediction is the seed ensemble
\begin{equation}
\bar p_c(x)=\frac15\sum_{k=0}^{4}\sigma\big(\ell^{(k)}_c(x)\big),
\label{eq:seedens}
\end{equation}
and all tests compare seed ensembles. Runs used Kaggle T4 GPUs; a run longer than one session resumed from an end-of-epoch checkpoint, which
gave the same predictions as uninterrupted training in a test.

\subsubsection{Encoder selection}
The plan screened four encoders on the NIH validation split (configuration \emph{orig}, seed 0, 10 epochs): DenseNet-121, ConvNeXt-V2-T
\cite{woo2023convnextv2}, RAD-DINO \cite{perezgarcia2025raddino} and Swin-V2-B \cite{liu2022swinv2}. The best encoder would go forward with
DenseNet-121, and an encoder within 0.005 of the best would yield to a cheaper one. All four reached validation macro-AUROC between 0.8384
and 0.8430, and DenseNet-121 was the cheapest, so it went forward alone. ConvNeXt-V2-T failed to train at the plan's learning rate;
following the rule of the suppressor study, the three encoders without batch normalisation were screened again with the encoder learning
rate reduced fourfold, a deviation decided before the other screen results were seen (Supplementary Section~S1).

\subsubsection{Finding groups}
Before any run, the classes were grouped as \emph{focal} (NIH: nodule, mass; VinBigData: nodule/mass), \emph{diffuse} (effusion,
infiltration, consolidation, oedema, pneumonia; VinBigData: pleural effusion, infiltration, consolidation, lung opacity, interstitial lung
disease) and \emph{boundary} (cardiomegaly, atelectasis, pneumothorax, pleural thickening, emphysema, fibrosis, hernia; VinBigData:
cardiomegaly, atelectasis, pneumothorax, pleural thickening, aortic enlargement, pulmonary fibrosis, calcification). VinBigData's ``other
lesion'' belongs to no group, and BDCXR-3257 (normal, bacterial or viral pneumonia) enters only the overall tests. For a group $G$ of classes
and two configurations $m$ and $m'$,
\begin{equation}
\begin{gathered}
\mathrm{AUROC}_{\mathrm{macro}}(m)=\frac1C\sum_{c=1}^{C}\mathrm{AUROC}_c(m),\\
\Delta_G(m,m')=\frac{1}{|G|}\sum_{c\in G}\big(\mathrm{AUROC}_c(m)-\mathrm{AUROC}_c(m')\big).
\end{gathered}
\label{eq:group}
\end{equation}

% ---------------------------------------------------------------------------------------------------------------------
\subsection{Study~III: Cross-Modality Knowledge Transfer from a 2.5D CT Teacher}
\label{sec:m-study3}
Studies~I and~II used CT only to render training images. Study~III asks whether CT can teach a radiograph classifier directly: a CT
classifier is distilled \cite{hinton2015distilling} into X-ray models on projections of the same volumes \cite{angaitkar2026anatomy}, and one
transformer family serves both modalities so that their heat maps can be compared in the same patient. Five hypotheses were fixed in a
fourth analysis plan, hash-frozen on 8 November 2026 before any model of this study was trained. Fig.~\ref{fig:xm_design} shows how the
data, the models and the evaluations of this study connect.

\begin{figure*}[!t]
\centering
\safefig[0.95\textwidth]{figures/fig_v4_design.pdf}
\caption{Study~III design. Data (left): the CT-RATE subset (\num{2156} volumes), the Qure.ai projections of CT-RATE volumes with six
shared findings (\num{1600} paired with the subset) and NIH ChestX-ray14 (14 diseases, 984 expert boxes), with VinBigData for zero-shot
transfer. Models (middle; SwinV2-T, five seeds each): the 2.5D CT teacher (24 slices in three Hounsfield windows), the projection models
$P_0$ and $P_1$ ($P_1$ distilled from the teacher's logits and pooled features) and the NIH classifiers $F_{\mathrm{in}}$,
$F_{\mathrm{p0}}$ and $F_{\mathrm{p1}}$, with DenseNet-121 as the convolutional reference. Evaluation (right): CT test AUROC and CT versus
X-ray heat maps (H14); projection test AUROC and zero-shot transfer; the NIH test split for H11--H13, the hallucination tests (H15), the
pointing game and the conformal false-negative rate.}
\label{fig:xm_design}
\end{figure*}

\subsubsection{Shared findings}
Six findings are shared between CT-RATE and the radiograph sets: lung opacity, pleural effusion, atelectasis, cardiomegaly, consolidation
and lung lesion. Each CT volume was reduced in the cloud environment to 24 axial slices over the lung extent with three Hounsfield windows
as channels (lung, mediastinum, bone), cropped to the body at 256\,px. NIH ChestX-ray14 and VinBigData are used with the splits of
Study~II. Pleural effusion, atelectasis, cardiomegaly and consolidation map to the classes of the same name and lung lesion to NIH Mass or
Nodule (VinBigData Nodule/Mass); these five are the exactly mapped shared findings. Lung opacity maps only approximately to NIH Infiltration
and is reported but not tested.

\subsubsection{Models and training}
All models use SwinV2-T (\texttt{swinv2\_tiny\_window8\_256}, ImageNet-1k weights \cite{liu2022swinv2,wightman2019timm}) with the
class-query decoder of Study~II, a learning rate of $3\times10^{-4}$ with the encoder at a quarter of it, model selection on the validation
split and five seeds per model (Fig.~\ref{fig:architecture}c). The \emph{CT teacher} is a 2.5D network: the encoder processes each of the
24 slices with the three windows as colour channels, each slice yields $8\times8$ tokens with a two-dimensional sine code and a learned
slice-depth embedding, and the decoder attends over the tokens of all slices (20 epochs, four volumes per step); a volumetric Swin network
did not fit a T4 GPU at a useful batch size. The \emph{projection model} $P_0$ learns the six findings from the projections (15 epochs,
batch 32); $P_1$ adds, on the \num{1600} training projections whose volume is in the CT subset, a distillation and an alignment term. For a
paired projection $k$ with student logits $z_k$, teacher logits $t_k$ and pooled decoder outputs $e_k$ (student) and $e^{\mathrm{T}}_k$
(teacher),
\begin{equation}
\begin{split}
\mathcal{L}={}&\mathcal{L}_{\mathrm{BCE}}(z,y)+\lambda_{\mathrm{kd}}\,T^{2}\,\mathcal{L}_{\mathrm{BCE}}\!\big(z_k/T,\,\sigma(t_k/T)\big)\\
&+\lambda_{\mathrm{al}}\big(1-\cos(We_k,\,e^{\mathrm{T}}_k)\big),
\end{split}
\label{eq:kd}
\end{equation}
with $y$ the six finding labels, $\mathcal{L}_{\mathrm{BCE}}$ the binary cross-entropy averaged over findings, $\sigma$ the logistic
function, $T=2$, $\lambda_{\mathrm{kd}}=1$, $\lambda_{\mathrm{al}}=0.5$ and $W$ a learned linear map; half of every $P_1$ batch is drawn
from the paired rows, and $P_1$ seed $k$ distils from the teacher of slot $k$. The \emph{radiograph classifiers} $F_{\mathrm{in}}$,
$F_{\mathrm{p0}}$ and $F_{\mathrm{p1}}$ start from ImageNet, $P_0$ and $P_1$ (encoder and decoder loaded, heads re-initialised) and are
trained on the 14 NIH diseases for 10 epochs with the radiograph as the only input. The five DenseNet-121 radiograph-only runs of Study~II
are the convolutional reference (trained for 20 epochs at the full encoder learning rate, the pre-specified rule for batch-normalised
encoders).

\subsubsection{Seed-replacement rule}
Teacher seed 2 did not learn (validation macro-AUROC 0.46--0.54 in every epoch, constant logits). Before any $P_1$ model of seed 2 was
trained we fixed a rule: train seeds 5 and 6 with the unchanged recipe and give slot 2 to seed 5 unless its logits were constant (per-class
standard deviation below 0.01), otherwise to seed 6. Seed 5 learned, so the teachers of slots 0--4 are seeds 0, 1, 5, 3 and 4
(Supplementary Table~S1).

% ---------------------------------------------------------------------------------------------------------------------
\subsection{Statistical Analysis Plan, Explainability and Hallucination Protocols}
\label{sec:m-stats}

All three studies share one statistical framework, stated first below; the hypotheses of each study and the protocols for explanations,
uncertainty and hallucination follow.

\subsubsection{Analysis plans and general statistics}
Each study was run under an analysis plan that was pre-specified and hash-frozen before the data it governs were evaluated, but not
deposited in a public registry: SAP.md (v1 family), SAP\_v2.md (Study~I, frozen on 1~November 2026; \path{freeze_v2.json} holds the hashes
of both plans, the SZCH development and held-out lists, the split files and the 52 weight files, at commit \texttt{dd52b6a}), the
BDCXR-3257 addendum SAP\_bdcxr.md, the Study~II plan (hashed on 3 November 2026), SAP\_v4.md with DESIGN\_v4.md (Study~III, frozen in
\path{freeze_v4.json} on 8 November 2026) and the exploratory site calibration SAP\_v5.md (frozen in \path{freeze_v5.json} on 17 November
2026). Every later change is logged in DEVIATIONS.md; all primary and secondary results are reported
whatever their direction.

All comparisons are paired. Confidence intervals are percentile bootstrap intervals \cite{efron1993bootstrap} (\num{10000} resamples unless
stated) over patients when several units come from one patient; paired differences are tested with the Wilcoxon signed-rank test
\cite{wilcoxon1945individual}; AUCs are compared with the paired DeLong test \cite{delong1988auc}. Seeds are averaged within unit before
testing and their spread is reported. Multiplicity is handled by the Holm procedure \cite{holm1979simple} within each family: with the $m$
$p$-values of a family ordered as $p_{(1)}\le\dots\le p_{(m)}$,
\begin{equation}
\tilde p_{(j)}=\max_{i\le j}\,\min\big\{1,\,(m-i+1)\,p_{(i)}\big\}.
\label{eq:holm}
\end{equation}
Confidence intervals are unadjusted 95\,\% intervals; effect sizes are always reported with them, and a hypothesis that fails is reported as
failed. Non-inferiority uses the paired mean difference $\bar\Delta=\frac1n\sum_i(m_i^{\mathrm{ours}}-m_i^{\mathrm{ref}})$ and the
one-sided 95\,\% lower bound $L_{95}$ of its image-level bootstrap distribution; it holds if
\begin{equation}
L_{95}(\bar\Delta)>-\delta
\label{eq:ni}
\end{equation}
for a margin $\delta$ fixed in the plan.

\subsubsection{Study~I hypotheses (H1$'$--H5$'$)}
\begin{description}
\item[H1$'$, dose-response.] On VinDr-RibCXR (245 images, six sites each, lesions from LIDC test CTs, $\alpha=1$), the per-site absorption
 $A$ averaged over seeds falls with the twin weight. With $\bar A_{s,w}$ the absorption at site $s$ averaged over the seeds trained with
 twin weight $w$, the slope $\beta_1$ of
 \begin{equation}
 \bar A_{s,w}=\beta_0+\beta_1\log_2(1+w)+e_{s,w}
 \label{eq:dose}
 \end{equation}
 has an image-cluster bootstrap interval, and H1$'$ passes if its upper bound is below zero. The same analysis on the Qure-projection
 test split (exact truth) is secondary.
\item[H2$'$, fidelity.] On the 51 held-out SZCH pairs, bone-region PSNR of ours against the reference model is non-inferior
 (Eq.~\eqref{eq:ni}) with margin $\delta=0.5$~dB; SSIM has a margin of 0.01 and is reported either way; ours against the no-twin control
 uses a margin of 0.25~dB (the cost of the regulariser). Metrics are computed calibrated (primary) and raw (Eq.~\eqref{eq:psnrcal}).
\item[H3$'$, lesion absorption.] On VinDr-RibCXR sites, $A$(ours)$<A$(reference model) and $A$(ours)$<A$(no-twin control), as paired
 site-level differences with image-cluster bootstrap intervals, Holm over the two contrasts, and an interaction between sites that overlap
 an annotated rib and sites that do not.
\item[H4$'$, evidence without insertion.] (a) $\Delta$CNR of real nodules on JSRT and Node21 (Eq.~\eqref{eq:cnr}), ours against the
 reference model and against the no-twin control, for all lesions and for lesions with at least 15\,\% bone overlap. (b) Downstream: NIH
 ChestX-ray14 macro-AUROC per input (original, ours, reference model, CLAHE, original plus ours; three seeds); TBX11K, Node21 and
 VinBigData detection by FROC and mAP@50; recall at two false positives per image for lesions with and without bone overlap, as a
 difference in differences against the original input.
\item[H5$'$, reliability.] Spearman correlation between ensemble spread and absolute bone error on the SZCH held-out pairs (pseudo-truth
 bone from the least-squares-aligned DES image); split-conformal coverage calibrated on Qure-projection validation CTs and measured on the
 Qure-projection test split (nominal 90\,\%) and, as measured under shift, on SZCH.
\end{description}
Site-level differences between models are paired, with image-cluster bootstrap intervals. Bootstrap intervals use \num{10000} resamples for
mean absorption and contrasts, \num{2000} for the overlap interaction and 500 for the dose-response slope, always resampling whole images
with replacement.

\emph{Exploratory site calibration (SAP\_v5, Section~\ref{sec:m-sitecal}).} H16 (primary) is a gain in the bone-region
PSNR$_{\mathrm{cal}}$ of ours from the calibration (one-sided 95\,\% lower bound of the paired difference above 0; image bootstrap,
\num{10000} resamples) and H16b the same for whole-image PSNR$_{\mathrm{cal}}$. H17 is the non-inferiority of the calibrated ours to the
calibrated reference model in bone-region PSNR$_{\mathrm{cal}}$ (margin 0.5\,dB, as H2$'$). H18 is a lower absorption of the calibrated ours
than of the calibrated reference model on VinDr-RibCXR, from one audit run with all four models so that they share every site
(image-cluster bootstrap, two-sided 95\,\% upper bound below 0). The rib benchmark of the calibrated models is descriptive. H16 is the only
confirmatory test of this plan; the others are reported without multiplicity correction.

\subsubsection{Study~II hypotheses (H6--H10)}
The hypotheses continue the numbering of the suppressor plans; H9 concerns the suppressor's intervals (Section~\ref{sec:m-ensemble}).
\begin{description}
\item[H6 (primary).] The hybrid has a higher macro-AUROC than \emph{orig}; one test per dataset, Holm correction over the three.
\item[H7 (secondary).] Hybrid against \emph{ctrl} (does a gain come from the suppression or from the extra computation?) and against
 \emph{early} (do separate streams beat channel stacking?); one Holm family over the six tests.
\item[H8 (secondary).] $\Delta_G$ of Eq.~\eqref{eq:group} for hybrid and for \emph{soft} against \emph{orig} in each finding group; Holm
 within dataset over the six tests. Per class, hybrid against \emph{orig} and against \emph{early} with the paired DeLong test
 \cite{delong1988auc}, Holm within dataset over classes and both contrasts.
\item[H10 (secondary).] In the hybrid models, the pointing-game hit rate of the soft-tissue stream minus that of the radiograph stream is
 larger for focal targets than for the other targets.
\end{description}
Every contrast $\Delta$ of two seed ensembles is resampled with a paired cluster bootstrap: replicate $b$ draws clusters with replacement
(patients for NIH, images for VinBigData, near-duplicate groups for BDCXR-3257), recomputes $\Delta^{*}_b$ with the same resamples for every
contrast of a dataset, and with $B=\num{10000}$ replicates gives the percentile interval $[\Delta^{*}_{(0.025)},\Delta^{*}_{(0.975)}]$ and
\begin{equation}
\begin{gathered}
p=\min\big\{1,\;2\min(p_{-},p_{+})\big\},\\
p_{-}=\frac{1+\#\{b:\Delta^{*}_b\le0\}}{B+1},\quad p_{+}=\frac{1+\#\{b:\Delta^{*}_b\ge0\}}{B+1}.
\end{gathered}
\label{eq:bootp}
\end{equation}
AUPRC differences are reported with bootstrap intervals only. The plan, the code and the split files were hashed on 3 November 2026, before
the encoder screen and before any classifier of this study saw a test split.

\subsubsection{Attribution maps and the pointing game}
\label{sec:m-xai}
For each hybrid model, Grad-CAM++ \cite{chattopadhay2018gradcampp} was computed per stream on the last encoder block. Integrated Gradients
\cite{sundararajan2017axiomatic} was computed on the three inputs jointly against a blurred baseline ($\sigma=8$\,px, 32 steps), because a
black image has no contrast after standardisation; the decoder attention was read from its last layer. A target $t$ is a pair of test image
and box class; a map $M_m$ hits when its maximum lies inside a box of that class, with a tolerance of 15\,px at 256\,px fixed in advance
\cite{zhang2018topdown}. With $h^{(k)}_m(t)\in\{0,1\}$ the hit of map $m$ of the seed-$k$ hybrid model, focal targets $\mathcal{F}$
(masses and nodules) and the other targets $\mathcal{O}$ (six box classes),
\begin{equation}
\begin{gathered}
d(t)=\frac15\sum_{k=0}^{4}\big(h^{(k)}_{S}(t)-h^{(k)}_{X}(t)\big),\\
I=\frac{1}{|\mathcal{F}|}\sum_{t\in\mathcal{F}}d(t)-\frac{1}{|\mathcal{O}|}\sum_{t\in\mathcal{O}}d(t),
\end{gathered}
\label{eq:pointing}
\end{equation}
and H10 tests $I>0$ with an image-cluster bootstrap (\num{10000} resamples; Eq.~\eqref{eq:bootp}).

\subsubsection{Classification uncertainty}
\label{sec:m-clsuq}
\emph{NIH: false-negative risk control.} An image's label set holds every class with ensemble probability at least $1-\lambda$; the threshold
is calibrated on the validation split (\num{8634} images) by conformal risk control \cite{angelopoulos2022crc} and the false-negative rate is
measured on the test split:
\begin{equation}
\begin{gathered}
\mathcal{C}_\lambda(x)=\{c:\bar p_c(x)\ge1-\lambda\},\\
\hat R_n(\lambda)=\frac1n\sum_{i=1}^{n}\Big(1-\frac{|\mathcal{C}_\lambda(x_i)\cap Y_i|}{|Y_i|}\Big),\\
\hat\lambda=\min\Big\{\lambda:\frac{n\,\hat R_n(\lambda)+1}{n+1}\le\alpha\Big\},
\end{gathered}
\label{eq:crc}
\end{equation}
over the $n$ calibration images with at least one positive label $Y_i$, at $\alpha=0.10$.

\emph{BDCXR-3257: class-conditional sets.} With probabilities normalised to sum to one, $\hat p_y=\bar p_y/\sum_{y'}\bar p_{y'}$, the score
$1-\hat p_y$ is calibrated separately for each true class \cite{vovk2005algorithmic}, on the other folds' out-of-fold predictions for each
test fold:
\begin{equation}
\begin{gathered}
\mathcal{C}(x)=\{y:1-\hat p_y(x)\le q_y\},\\
q_y=\text{the }j_y\text{-th smallest of }\{1-\hat p_{y_i}(x_i):y_i=y\},\\
j_y=\lceil (n_y+1)(1-\alpha)\rceil.
\end{gathered}
\label{eq:mondriansets}
\end{equation}

\emph{Selective classification.} Images are ranked by a confidence score and the most confident fraction is kept
\cite{geifman2017selective,geifman2019bias}. Score A is one minus the entropy $H$ of the seed ensemble (mean binary entropy over the 14
labels for NIH, Shannon entropy of the normalised probabilities for BDCXR-3257); score B adds the suppressor's own uncertainty,
$-(z_H+z_\sigma)$, where $\sigma$ is the spread of Eq.~\eqref{eq:ensemble} averaged over the lung field and $z$ standardises over the
evaluated images of a dataset. With the losses $\ell_{(1)},\dots,\ell_{(N)}$ sorted by decreasing confidence (Brier score for NIH, top-1
error for BDCXR-3257),
\begin{equation}
\mathrm{AURC}=\frac1N\sum_{k=1}^{N}\frac1k\sum_{j=1}^{k}\ell_{(j)}.
\label{eq:aurc}
\end{equation}

\subsubsection{Study~III: cross-modal agreement and hallucination tests}
\label{sec:m-xmxai}
For the radiograph classifiers we computed Grad-CAM++ on the last encoder stage \cite{chattopadhay2018gradcampp}, Integrated Gradients
\cite{sundararajan2017axiomatic} and the attention of the last decoder layer, scored with the pointing game on the 984 NIH expert boxes
\cite{zhang2018topdown} at the 15\,px tolerance of Study~II. For each CT test scan positive for a finding, the teacher's slice-wise
Grad-CAM++ maps were stacked and summed along the anteroposterior axis, giving a frontal map in the geometry of the paired projection. CT and
X-ray maps were cropped to their own lung boxes, resized to $32\times32$, scaled to sum to one and averaged over the five seed pairs.
Agreement is the Pearson correlation of the matched maps, and the null correlates the CT map with the X-ray maps of the other scans positive
for the same finding. A pre-specified orientation check on the attenuation frames alone (correlation 0.472 as they were, 0.435 mirrored)
left the maps unmirrored.

We call a report of a finding that the image does not contain a hallucination. On the \num{9861} NIH test images without any finding, with
seed-mean probabilities and per-class thresholds at 90\,\% specificity on the validation split, the hallucination rate is the share of
images in which some class moves from at most to above its threshold after a perturbation that adds no disease: rib emphasis (the
radiograph plus half of the bone image of the Study~I suppressor), Gaussian noise ($\sigma=0.02$) or a gamma of 0.8. Counterfactual nodules
from the audit library (Section~\ref{sec:m-audit}) were pasted into lung fields of images whose Mass and Nodule scores were at most their
thresholds; sensitivity is the share in which either score crosses it. Off-target attribution is the share of Grad-CAM++ mass outside the
lung fields for positive predictions of lung findings.

\subsubsection{Study~III hypotheses (H11--H15)}
A model's prediction is the mean of its seeds' probabilities. H11 (primary) compares $F_{\mathrm{p1}}$ with $F_{\mathrm{in}}$ in the mean
NIH test AUROC of the five shared findings; H12 compares $F_{\mathrm{p1}}$ with $F_{\mathrm{p0}}$ and $F_{\mathrm{p0}}$ with
$F_{\mathrm{in}}$; H13 compares $F_{\mathrm{in}}$ with DenseNet-121 in the 14-disease macro-AUROC. That reference had been evaluated on the
NIH test split in Study~II about ten hours before this plan was frozen, so its score was known when H13 was specified (no Swin model existed
yet). These contrasts use a paired patient-cluster bootstrap (\num{10000} resamples, two-sided percentile $p$-values, Eq.~\eqref{eq:bootp});
H11 and the two H12 contrasts form one Holm family. H14 is the mean per-map difference between matched and mismatched correlation
(scan-cluster bootstrap) and H15 the mean per-image difference in the rib-emphasis flip between $F_{\mathrm{p1}}$ and $F_{\mathrm{in}}$
(patient-cluster bootstrap). The CT teacher, the projection models and their zero-shot transfer, per-class DeLong tests
\cite{delong1988auc}, conformal risk control of the false-negative rate at 0.10 \cite{angelopoulos2022crc} and the pointing game are
descriptive; the grouping of the DeLong tests and the pointing-game summary were fixed after the hypothesis results were known.
